\documentclass{article}
\usepackage{amsmath}
% a comment with $not math$ and \begin{equation} nothing
\newcommand{\R}{\mathbb{R}}
\newcommand{\norm}[1]{\left\lVert #1\right\rVert}
\newcommand{\inner}[2][x]{\langle #1, #2\rangle}
\renewcommand{\epsilon}{\varepsilon}
\def\E{\mathbb{E}}
\def\pair#1#2{(#1,#2)}
\DeclareMathOperator{\Var}{Var}
\DeclareMathOperator*{\argmin}{arg\,min}
\let\ip\inner
\newenvironment{eq}{\begin{equation}}{\end{equation}}
\newcommand{\be}{\begin{equation}}
\newcommand{\ee}{\end{equation}}
\newcommand{\twice}[1]{#1#1}
\newcommand{\loopy}{\loopy x}
\begin{document}
\title{Fixture $\alpha$}
\maketitle
\section{Introduction}
Let $x \in \R^n$ with $\norm{x} \le 1$ and $\Var(X)=\E[X^2]$.
The price is \$5 and the sum is $a+b$.
More text with \(y = 2x\) inline, and a display
\[
  \int_0^1 g(t)\,dt = 1 .
\]
$$ a^2 + b^2 = c^2 $$

\begin{equation}\label{eq:one}
  f(x) = \sum_{i=1}^{n} \inner[y]{x}, \qquad \pair{a}{b}.
\end{equation}
\be
\argmin_x \twice{y}
\ee
\begin{eq}
\epsilon = 0
\end{eq}
\end{document}
